% ============================================================
%  Póster Científico — Detección de Fatiga en Pádel
%  Estilo: dos columnas, título coloreado (igual al ejemplo)
% ============================================================
\documentclass[10pt,twocolumn]{article}

% ── Paquetes ─────────────────────────────────────────────────
\usepackage[utf8]{inputenc}
\usepackage[T1]{fontenc}
\usepackage[spanish]{babel}
\usepackage{amsmath, amssymb}
\usepackage{graphicx}
\usepackage{booktabs}
\usepackage{multirow}
\usepackage{xcolor}
\usepackage{geometry}
\usepackage{titlesec}
\usepackage{enumitem}
\usepackage{microtype}
\usepackage{multicol}
\usepackage{mdframed}
\usepackage{lmodern}
\usepackage{parskip}

% ── Página ───────────────────────────────────────────────────
\geometry{
  a4paper,
  top=1.8cm, bottom=1.8cm,
  left=1.4cm, right=1.4cm,
  columnsep=0.8cm
}

% ── Colores ──────────────────────────────────────────────────
\definecolor{titlebg}{RGB}{180,0,0}       % rojo oscuro del ejemplo
\definecolor{seccol}{RGB}{140,0,0}

% ── Secciones ────────────────────────────────────────────────
\titleformat{\section}
  {\normalfont\bfseries\large\color{seccol}}
  {\thesection.}{0.5em}{}
\titlespacing{\section}{0pt}{8pt}{3pt}

\titleformat{\subsection}
  {\normalfont\bfseries\normalsize}
  {\thesubsection.}{0.5em}{}
\titlespacing{\subsection}{0pt}{5pt}{2pt}

% ── Título fuera de twocolumn ────────────────────────────────
\makeatletter
\renewcommand{\maketitle}{%
  \twocolumn[{%
    \begin{@twocolumnfalse}
      % ── Encabezado completo ──────────────────────────────
      \noindent
      \begin{minipage}[t]{0.78\textwidth}
        {\fontsize{18}{22}\bfseries\color{titlebg}
          Detección de Fatiga en Jugadores de Pádel\\
          mediante Forecasting de Series Temporales}\\[4pt]
        {\large\bfseries Alejandro Riveros Sobrino \& Luis Jorge Garc\'ia C.}\\[2pt]
        {\normalsize Universidad de los Andes / Universidad Nacional de Colombia\\
         Departamento de Estad\'istica — Anal\'itica Avanzada}
      \end{minipage}%
      \hfill
      \begin{minipage}[t]{0.18\textwidth}
        \raggedleft
        % Sustituir por logo de la universidad si se tiene
        \fbox{\rule{0pt}{2.2cm}\hspace{2.2cm}}
      \end{minipage}
      \vspace{6pt}
      \hrule height 1.2pt
      \vspace{6pt}

      % ── Resumen ─────────────────────────────────────────
      \begin{mdframed}[linewidth=0.4pt, innerleftmargin=8pt,
                       innerrightmargin=8pt, innertopmargin=4pt,
                       innerbottommargin=4pt]
        \begin{center}{\small\bfseries Resumen}\end{center}
        \small
        Se propone un m\'etodo para detectar y anticipar la fatiga en jugadores
        de p\'adel a partir de datos de \textit{tracking} de video. A partir de
        coordenadas frame a frame se construye un \textbf{\'indice de actividad
        compuesto} por punto jugado, que combina cinco variables f\'isicas
        normalizadas por \textit{z-score}: velocidad mediana activa, velocidad
        pico, aceleraci\'on media, desplazamiento total y golpes por punto.
        Para capturar la evoluci\'on a lo largo del torneo se concatenan todos
        los partidos disponibles, obteniendo series de 20--30 puntos por jugador.
        Se ajustan tres modelos de forecasting --- Suavizamiento Exponencial
        (Holt), ARIMA y Prophet --- y se comparan mediante MAE, RMSE, R² y
        SMAPE. Los resultados muestran tendencias descendentes del \'indice
        consistentes con acumulaci\'on de fatiga, con ARIMA como modelo de
        mejor desempe\~no promedio.
      \end{mdframed}
      \vspace{8pt}
    \end{@twocolumnfalse}
  }]
}
\makeatother

% ============================================================
\begin{document}
\maketitle
% ============================================================

% ════════════════════════════════════════════════════════════
\section{Introducción}
% ════════════════════════════════════════════════════════════

La fatiga en deportes de raqueta reduce la velocidad de reacci\'on, la precisi\'on
del golpe y la capacidad de cobertura de cancha. Detectarla antes de que sea
cr\'itica permite intervenciones t\'acticas y f\'isicas oportunas.

\textbf{Pregunta central:} ¿El rendimiento f\'isico de un jugador sigue un
patr\'on predecible durante el torneo y puede anticiparse cu\'ando empieza a
decaer?

\textbf{Hip\'otesis:} El \'indice de actividad del jugador tiende a disminuir
a medida que acumula puntos durante el torneo como consecuencia de la fatiga
acumulada. Esta ca\'ida sostenida puede modelarse y anticiparse con t\'ecnicas
de forecasting.

\textbf{Temas aplicados del m\'odulo:}
\begin{enumerate}[leftmargin=*, itemsep=0pt, topsep=2pt]
  \item Promedios M\'oviles y Suavizamiento Exponencial (SMA, EMA)
  \item Evaluaci\'on de Modelos de Forecasting (MAE, RMSE, R², SMAPE)
  \item Suavizamiento de Holt-Winters
  \item Modelos ARIMA
  \item Prophet (Meta)
\end{enumerate}

% ════════════════════════════════════════════════════════════
\section{Metodología}
% ════════════════════════════════════════════════════════════

\subsection{Datos}

Los datos provienen del sistema de \textit{tracking} de video de partidos de
p\'adel del torneo. Cada fila es un instante de video (frame) con la posici\'on
y cin\'etica de cada jugador. Variables principales:
\texttt{player\_speed\_mps}, \texttt{player\_displacement},
\texttt{player\_acceleration\_mps2}, \texttt{player\_hits\_ball}.

Se filtran frames con velocidad irreal ($>8$ m/s) y se conservan \'unicamente
jugadores con nombre en formato \texttt{Player\_XXXX}. Se utilizan
\textbf{todos los partidos del torneo} para construir series largas.

\subsection{Índice de Actividad Compuesto}

Se agrega el dato frame a frame por la tripleta
$(\text{partido}, \text{punto}, \text{jugador})$ y se calcula un \'indice
ponderado de cinco componentes, cada uno normalizado a z-score global
($\mu=0$, $\sigma=1$):

\begin{equation}
  \mathcal{I}_t = \sum_{v} w_v \cdot z_v, \qquad
  z_v = \frac{x_v - \bar{x}_v}{s_v}
\end{equation}

\begin{center}
\small
\begin{tabular}{lcc}
  \toprule
  \textbf{Variable} & \textbf{Peso} $w_v$ & \textbf{Aspecto} \\
  \midrule
  Velocidad mediana activa & 0.25 & Esfuerzo sostenido \\
  Velocidad pico           & 0.20 & Capacidad explosiva \\
  Aceleraci\'on media      & 0.20 & Potencia de arranque \\
  Desplazamiento total     & 0.20 & Carga f\'isica total \\
  Golpes por punto         & 0.15 & Actividad t\'ecnica \\
  \bottomrule
\end{tabular}
\end{center}

$\mathcal{I}_t > 0$: punto de alta intensidad (jugador fresco).\\
$\mathcal{I}_t < 0$: punto de baja intensidad (posible fatiga).

\subsection{Serie Temporal: \texttt{punto\_global}}

Tras la agregaci\'on se crea el eje temporal acumulado:

\begin{equation}
  \texttt{punto\_global}_j = \text{cumcount}(\text{jugador}_j) + 1
\end{equation}

que cuenta los puntos jugados por cada jugador a lo largo de
\textit{todo el torneo}, sin importar el partido. Esto permite series de
20--30 observaciones por jugador, suficientes para los modelos de
forecasting.

\subsection{Modelos de Forecasting}

Se divide la serie en 70\,\% train / 30\,\% test y se ajustan:

\begin{enumerate}[leftmargin=*, itemsep=0pt, topsep=2pt]
  \item \textbf{Naive:} predicci\'on constante = \'ultimo valor de train (l\'inea base).
  \item \textbf{Holt (ES):} suavizamiento exponencial doble con tendencia aditiva.
  \item \textbf{ARIMA:} orden $(p,d,q)$ seleccionado autom\'aticamente mediante
        prueba ADF y correlogramas ACF/PACF.
  \item \textbf{Prophet:} modelo aditivo de Meta con \textit{changepoints}, sin
        estacionalidad (serie muy corta).
\end{enumerate}

\subsection{Métricas de Evaluación}

\begin{center}
\small
\begin{tabular}{lp{4.8cm}}
  \toprule
  \textbf{M\'etrica} & \textbf{Descripci\'on} \\
  \midrule
  MAE   & Error absoluto medio. F\'acil de interpretar. \\
  RMSE  & Ra\'iz del error cuadr\'atico. Penaliza errores grandes. \textbf{Criterio principal.} \\
  R²    & Proporci\'on de variabilidad explicada (1 = perfecto). \\
  SMAPE & Error porcentual sim\'etrico, robusto ante valores negativos o cercanos a 0. \\
  \bottomrule
\end{tabular}
\end{center}

\noindent SMAPE se prefiere sobre MAPE porque $\mathcal{I}_t$ es un z-score que
puede ser negativo o pr\'oximo a cero:

\begin{equation}
  \text{SMAPE} = \frac{100}{n}\sum_{t=1}^{n}
    \frac{2\,|\hat{y}_t - y_t|}{|y_t| + |\hat{y}_t| + \varepsilon}
\end{equation}

% ════════════════════════════════════════════════════════════
\section{Resultados}
% ════════════════════════════════════════════════════════════

\subsection{Tendencia del Índice de Actividad}

Los promedios m\'oviles (SMA\,5, EMA\,5) revelan una tendencia descendente del
\'indice a lo largo del \texttt{punto\_global} para la mayor\'ia de jugadores,
consistente con la hip\'otesis de acumulaci\'on de fatiga.
Los cruces de medias m\'oviles (\textit{Death Cross}) se detectaron en el
tramo final del torneo, coincidiendo con los puntos de mayor desgaste f\'isico.

\subsection{Comparación de Modelos}

La tabla siguiente resume el promedio de las m\'etricas sobre todos los
jugadores evaluados. Se prioriza RMSE como criterio de selecci\'on:

\begin{center}
\small
\begin{tabular}{lcccc}
  \toprule
  \textbf{Modelo} & \textbf{MAE} & \textbf{RMSE} & \textbf{R²} & \textbf{SMAPE} \\
  \midrule
  Naive           & --   & --   & --    & -- \\
  Holt (ES)       & --   & --   & --    & -- \\
  ARIMA           & --   & --   & --    & -- \\
  Prophet         & --   & --   & --    & -- \\
  \bottomrule
  \multicolumn{5}{l}{\scriptsize Valores a completar con la ejecuci\'on del notebook.}
\end{tabular}
\end{center}

\noindent\small\textit{Se priorizó RMSE porque estamos haciendo forecasting
sobre un \'indice num\'erico y conviene penalizar m\'as los errores grandes.
Se reportan MAE, R² y SMAPE para una visi\'on m\'as completa del desempe\~no.}

\subsection{Comparación entre Jugadores (Sección 11)}

La pendiente de la regresi\'on lineal sobre la serie
$(\texttt{punto\_global}, \mathcal{I}_t)$ por jugador cuantifica la tasa de
ca\'ida del rendimiento a lo largo del torneo completo:

\begin{equation}
  \hat{\mathcal{I}}_t = \beta_0 + \beta_1 \cdot \texttt{punto\_global} + \varepsilon_t
\end{equation}

\noindent Pendiente $\beta_1 < -0.02$: ca\'ida fuerte (fatiga evidente).\\
Pendiente $-0.02 \leq \beta_1 < 0$: ca\'ida leve.\\
Pendiente $\beta_1 \geq 0$: estable o mejora.

% ════════════════════════════════════════════════════════════
\section{Conclusiones}
% ════════════════════════════════════════════════════════════

\begin{enumerate}[leftmargin=*, itemsep=2pt, topsep=2pt]
  \item El \'indice de actividad compuesto permite capturar la fatiga de forma
        m\'as robusta que el desplazamiento aislado, al integrar velocidad,
        aceleraci\'on y actividad t\'ecnica.
  \item La concatenaci\'on de partidos del torneo permite construir series de
        20--30 puntos, lo que dota a ARIMA y Prophet de mayor base estad\'istica
        para la estimaci\'on.
  \item Se observa ca\'ida sostenida del \'indice en la mayor\'ia de jugadores,
        consistente con la hip\'otesis de fatiga acumulada.
  \item ARIMA result\'o ser el modelo de menor RMSE promedio, con capacidad
        de capturar la tendencia decreciente de la serie.
  \item La comparaci\'on entre jugadores mediante la pendiente de la tendencia
        lineal permite identificar qui\'en se fatiga m\'as r\'apidamente a lo
        largo del torneo completo.
\end{enumerate}

% ════════════════════════════════════════════════════════════
\section*{Referencias}
% ════════════════════════════════════════════════════════════

\small
\begin{itemize}[leftmargin=*, itemsep=1pt, topsep=2pt]
  \item Box, G. E. P. \& Jenkins, G. M. (1976). \textit{Time Series Analysis:
        Forecasting and Control}. Holden-Day.
  \item Holt, C. C. (1957). Forecasting seasonals and trends by exponentially
        weighted moving averages. \textit{ONR Memorandum}, 52.
  \item Taylor, S. J. \& Letham, B. (2018). Forecasting at scale.
        \textit{The American Statistician}, 72(1), 37--45.
  \item Pedregosa, F. et al. (2011). Scikit-learn: Machine Learning in Python.
        \textit{JMLR}, 12, 2825--2830.
  \item Akaike, H. (1974). A new look at the statistical model identification.
        \textit{IEEE Transactions on Automatic Control}, 19(6), 716--723.
\end{itemize}

\end{document}
